\begin{abstract}
Laboratory midge swarms remain cohesive despite weak global velocity
alignment, providing a system for testing whether local organization persists
after ordinary neighborhood deformation is removed. We re-analysed 19
three-dimensional \textit{Chironomus riparius} swarm recordings and defined T1
as tangential non-affine activity remaining after finite-lag local affine motion
was fitted and subtracted. Under the frozen centered-detrending pipeline and a
two-scale consistency criterion, T1 was supported in 14 of 19 observations.
None of 1000 full-pipeline pseudo-event null replicates reached this count
(empirical \(p \approx 0.001\)). Support remained at majority level under
past-only detrending (11/19) and without rolling detrending (13/19), showing
that the phenomenon is common but the exact count is preprocessing-dependent.
Within supported observations, the result was stable to nearby neighborhood
scales and temporal lags, and diffuse tangential activity was the most
consistent phenotype. Compact-state variables did not improve held-out first-
or second-moment prediction consistently, true transition times showed no
robust near-pre excess after state matching, and recent state-path history was
informative only in some observations. T1 therefore defines a measurable local
residual layer rather than identifying a universal mechanism.
\end{abstract}

\begin{IEEEkeywords}
midge swarms, non-affine motion, local affine subtraction, collective motion,
compact state variables, observation heterogeneity
\end{IEEEkeywords}
